% --- Archon physics formalization source begin ---
% archon:physics
% archon:covers PhyXMiniProblems/problem_phyx_mini_0497.lean
% archon:source-report reports/phyx_mini/problem_phyx_mini_0497.source.json

\chapter{Physics problem phyx\_mini\_0497}
\label{ch:PhyXMiniProblems_problem_phyx_mini_0497}

\paragraph{Problem source.}
\#\# Physical scenario

An alpha particle approaches a \textbackslash{}( \textasciicircum{}\{197\}\textbackslash{}text\{Au\} \textbackslash{}) nucleus with a speed of \textbackslash{}( 1.50 \textbackslash{}times 10\textasciicircum{}7 \textbackslash{}text\{ m/s\} \textbackslash{}). As figure shows, the alpha particle is scattered at a \textbackslash{}( 49\textasciicircum{}\textbackslash{}circ \textbackslash{}) angle at the slower speed of \textbackslash{}( 1.49 \textbackslash{}times 10\textasciicircum{}7 \textbackslash{}text\{ m/s\} \textbackslash{}).

\#\# Question

In what speed does the \textbackslash{}( \textasciicircum{}\{197\}\textbackslash{}text\{Au\} \textbackslash{}) nucleus recoil?

\#\# Answer choices

- A: A: \textbackslash{}( 4.08 \textbackslash{}times 10\textasciicircum{}5 \textbackslash{}text\{ m/s\} \textbackslash{})

- B: B: \textbackslash{}( 1.26 \textbackslash{}times 10\textasciicircum{}5 \textbackslash{}text\{ m/s\} \textbackslash{})

- C: C: \textbackslash{}( 2.52 \textbackslash{}times 10\textasciicircum{}5 \textbackslash{}text\{ m/s\} \textbackslash{})

- D: D: \textbackslash{}( 5.04 \textbackslash{}times 10\textasciicircum{}5 \textbackslash{}text\{ m/s\} \textbackslash{})

\#\# Figure caption (auxiliary; use the image as primary evidence)

The image illustrates a diagram of an alpha particle scattering experiment. Here's a breakdown of the content:

- There are two main objects: an alpha particle (α) and a gold (Au) nucleus.
- The alpha particle is represented by a small circle with plus signs, indicating positive charge.
- The gold nucleus is also depicted by a larger circle with plus signs, signifying more positive charge.
- A blue curved arrow shows the path of the alpha particle as it approaches and then is deflected by the gold nucleus.
- The angle of deflection is indicated as 49° between a horizontal line extending from the nucleus and the path of the deflected alpha particle.
- A green arrow points in the direction of the deflection.
- The text includes "α" labeling the alpha particle and "197 Au nucleus" identifying the gold nucleus, with the superscript "197" representing the isotope number of gold.

The diagram illustrates the concept of nuclear scattering, particularly how alpha particles are deflected by positively charged gold nuclei.

\#\# Dataset metadata

Nuclear Physics; Physical Model Grounding Reasoning, Multi-Formula Reasoning

\paragraph{Recorded answer/context.}
C: C: \textbackslash{}( 2.52 \textbackslash{}times 10\textasciicircum{}5 \textbackslash{}text\{ m/s\} \textbackslash{})

\paragraph{Figure/image path.}
/root/proposal\_for\_physic/science-mango/phyx\_mini\_run/phyx\_data/test\_image/497.png

\paragraph{Formalization target.}
create a compiling Lean file with sorry bodies at `PhyXMiniProblems/problem\_phyx\_mini\_0497.lean`. The Lean declarations must preserve the physical quantities, dimensions or dimensional roles, figure labels, governing-law hypotheses, and final relation expressed by this problem.
Use Mathlib/Physlib names found through LeanExplore where available. If a physics API is missing, introduce faithful local abstractions rather than scalar placeholder aliases.

\begin{theorem}[Physics formalization target]
\label{thm:physics:phyx_mini_0497:target}
\lean{PhyXMiniProblems.ProblemPhyXMini0497.problem_phyx_mini_0497}
\uses{def:physics:phyx-mini-0497:phyxminiproblems-problemphyxmini0497-speedinmeterspersecond, def:physics:phyx-mini-0497:phyxminiproblems-problemphyxmini0497-degreestoradians, def:physics:phyx-mini-0497:phyxminiproblems-problemphyxmini0497-alphagoldscatteringsetup, def:physics:phyx-mini-0497:phyxminiproblems-problemphyxmini0497-incomingalphaspeeddatum, def:physics:phyx-mini-0497:phyxminiproblems-problemphyxmini0497-scatteredalphaspeeddatum, def:physics:phyx-mini-0497:phyxminiproblems-problemphyxmini0497-scatteringangledegreesdatum, def:physics:phyx-mini-0497:phyxminiproblems-problemphyxmini0497-matchesproblemandprimaryfigure, def:physics:phyx-mini-0497:phyxminiproblems-problemphyxmini0497-satisfiesmassnumberapproximation, def:physics:phyx-mini-0497:phyxminiproblems-problemphyxmini0497-satisfiesscatteringgeometry, def:physics:phyx-mini-0497:phyxminiproblems-problemphyxmini0497-satisfiesmomentumconservation, def:physics:phyx-mini-0497:phyxminiproblems-problemphyxmini0497-recordedanswerchoice, def:physics:phyx-mini-0497:phyxminiproblems-problemphyxmini0497-matchesanswerchoice}
The assigned autoformalize agent should translate this physics problem into Lean declarations in the covered file, with theorem and lemma proof bodies written as `by sorry`.
\end{theorem}
\begin{proof}
This is an autoformalization task, not a proof task. Produce statements that can later be proved by the physics prover without weakening the physical model.
\end{proof}

% --- Archon declaration topology begin ---
\section{Declaration topology}
\begin{definition}[Mass Quantity]
  \label{def:physics:phyx-mini-0497:phyxminiproblems-problemphyxmini0497-massquantity}
  \lean{PhyXMiniProblems.ProblemPhyXMini0497.MassQuantity}
  A nonnegative physical mass, independent of the unit used to read it
\end{definition}
\begin{proof}
  This is the declared model definition of Mass Quantity.
\end{proof}
\begin{definition}[Planar Axis]
  \label{def:physics:phyx-mini-0497:phyxminiproblems-problemphyxmini0497-planaraxis}
  \lean{PhyXMiniProblems.ProblemPhyXMini0497.PlanarAxis}
  The horizontal and vertical axes in the scattering plane
\end{definition}
\begin{proof}
  This is the declared model definition of Planar Axis.
\end{proof}
\begin{definition}[Planar Velocity]
  \label{def:physics:phyx-mini-0497:phyxminiproblems-problemphyxmini0497-planarvelocity}
  \lean{PhyXMiniProblems.ProblemPhyXMini0497.PlanarVelocity}
  \uses{def:physics:phyx-mini-0497:phyxminiproblems-problemphyxmini0497-planaraxis}
  A signed two-dimensional physical velocity with dimension length per time
\end{definition}
\begin{proof}
  This is the declared model definition of Planar Velocity.
\end{proof}
\begin{definition}[mass Readout]
  \label{def:physics:phyx-mini-0497:phyxminiproblems-problemphyxmini0497-massreadout}
  \lean{PhyXMiniProblems.ProblemPhyXMini0497.massReadout}
  \uses{def:physics:phyx-mini-0497:phyxminiproblems-problemphyxmini0497-massquantity}
  Read a physical mass in a selected mass unit
\end{definition}
\begin{proof}
  This is the declared model definition of mass Readout.
\end{proof}
\begin{definition}[velocity Readout]
  \label{def:physics:phyx-mini-0497:phyxminiproblems-problemphyxmini0497-velocityreadout}
  \lean{PhyXMiniProblems.ProblemPhyXMini0497.velocityReadout}
  \uses{def:physics:phyx-mini-0497:phyxminiproblems-problemphyxmini0497-planaraxis, def:physics:phyx-mini-0497:phyxminiproblems-problemphyxmini0497-planarvelocity}
  Read a planar velocity in coherent selected length and time units
\end{definition}
\begin{proof}
  This is the declared model definition of velocity Readout.
\end{proof}
\begin{definition}[speed In Meters Per Second]
  \label{def:physics:phyx-mini-0497:phyxminiproblems-problemphyxmini0497-speedinmeterspersecond}
  \lean{PhyXMiniProblems.ProblemPhyXMini0497.speedInMetersPerSecond}
  \uses{def:physics:phyx-mini-0497:phyxminiproblems-problemphyxmini0497-planarvelocity, def:physics:phyx-mini-0497:phyxminiproblems-problemphyxmini0497-velocityreadout}
  The Euclidean speed readout in metres per SI second
\end{definition}
\begin{proof}
  This is the declared model definition of speed In Meters Per Second.
\end{proof}
\begin{definition}[degrees To Radians]
  \label{def:physics:phyx-mini-0497:phyxminiproblems-problemphyxmini0497-degreestoradians}
  \lean{PhyXMiniProblems.ProblemPhyXMini0497.degreesToRadians}
  Convert degrees to the radian argument used by Real.sin and Real.cos
\end{definition}
\begin{proof}
  This is the declared model definition of degrees To Radians.
\end{proof}
\begin{definition}[Figure Object]
  \label{def:physics:phyx-mini-0497:phyxminiproblems-problemphyxmini0497-figureobject}
  \lean{PhyXMiniProblems.ProblemPhyXMini0497.FigureObject}
  The two physical bodies distinguished in the scattering diagram
\end{definition}
\begin{proof}
  This is the declared model definition of Figure Object.
\end{proof}
\begin{definition}[Figure Label]
  \label{def:physics:phyx-mini-0497:phyxminiproblems-problemphyxmini0497-figurelabel}
  \lean{PhyXMiniProblems.ProblemPhyXMini0497.FigureLabel}
  Text and numerical labels visible in the primary image
\end{definition}
\begin{proof}
  This is the declared model definition of Figure Label.
\end{proof}
\begin{definition}[Velocity Arrow]
  \label{def:physics:phyx-mini-0497:phyxminiproblems-problemphyxmini0497-velocityarrow}
  \lean{PhyXMiniProblems.ProblemPhyXMini0497.VelocityArrow}
  The three velocity arrows relevant to momentum accounting
\end{definition}
\begin{proof}
  This is the declared model definition of Velocity Arrow.
\end{proof}
\begin{definition}[Alpha Gold Scattering Figure]
  \label{def:physics:phyx-mini-0497:phyxminiproblems-problemphyxmini0497-alphagoldscatteringfigure}
  \lean{PhyXMiniProblems.ProblemPhyXMini0497.AlphaGoldScatteringFigure}
  \uses{def:physics:phyx-mini-0497:phyxminiproblems-problemphyxmini0497-figureobject, def:physics:phyx-mini-0497:phyxminiproblems-problemphyxmini0497-figurelabel, def:physics:phyx-mini-0497:phyxminiproblems-problemphyxmini0497-velocityarrow}
  Qualitative facts that may be read directly from the supplied diagram
\end{definition}
\begin{proof}
  This is the declared model definition of Alpha Gold Scattering Figure.
\end{proof}
\begin{definition}[Alpha Gold Scattering Setup]
  \label{def:physics:phyx-mini-0497:phyxminiproblems-problemphyxmini0497-alphagoldscatteringsetup}
  \lean{PhyXMiniProblems.ProblemPhyXMini0497.AlphaGoldScatteringSetup}
  \uses{def:physics:phyx-mini-0497:phyxminiproblems-problemphyxmini0497-massquantity, def:physics:phyx-mini-0497:phyxminiproblems-problemphyxmini0497-planarvelocity, def:physics:phyx-mini-0497:phyxminiproblems-problemphyxmini0497-alphagoldscatteringfigure}
  All physical quantities in the two-body event. The gold recoil velocity is an unconstrained dimensionful vector here; its requested magnitude is not a field of the setup
\end{definition}
\begin{proof}
  This is the declared model definition of Alpha Gold Scattering Setup.
\end{proof}
\begin{definition}[incoming Alpha Speed Datum]
  \label{def:physics:phyx-mini-0497:phyxminiproblems-problemphyxmini0497-incomingalphaspeeddatum}
  \lean{PhyXMiniProblems.ProblemPhyXMini0497.incomingAlphaSpeedDatum}
  The incident alpha speed stated in the problem, in metres per second
\end{definition}
\begin{proof}
  This is the declared model definition of incoming Alpha Speed Datum.
\end{proof}
\begin{definition}[scattered Alpha Speed Datum]
  \label{def:physics:phyx-mini-0497:phyxminiproblems-problemphyxmini0497-scatteredalphaspeeddatum}
  \lean{PhyXMiniProblems.ProblemPhyXMini0497.scatteredAlphaSpeedDatum}
  The scattered alpha speed stated in the problem, in metres per second
\end{definition}
\begin{proof}
  This is the declared model definition of scattered Alpha Speed Datum.
\end{proof}
\begin{definition}[scattering Angle Degrees Datum]
  \label{def:physics:phyx-mini-0497:phyxminiproblems-problemphyxmini0497-scatteringangledegreesdatum}
  \lean{PhyXMiniProblems.ProblemPhyXMini0497.scatteringAngleDegreesDatum}
  The alpha-particle scattering angle shown in the figure, in degrees
\end{definition}
\begin{proof}
  This is the declared model definition of scattering Angle Degrees Datum.
\end{proof}
\begin{definition}[Matches Problem And Primary Figure]
  \label{def:physics:phyx-mini-0497:phyxminiproblems-problemphyxmini0497-matchesproblemandprimaryfigure}
  \lean{PhyXMiniProblems.ProblemPhyXMini0497.MatchesProblemAndPrimaryFigure}
  \uses{def:physics:phyx-mini-0497:phyxminiproblems-problemphyxmini0497-velocityreadout, def:physics:phyx-mini-0497:phyxminiproblems-problemphyxmini0497-speedinmeterspersecond, def:physics:phyx-mini-0497:phyxminiproblems-problemphyxmini0497-alphagoldscatteringsetup, def:physics:phyx-mini-0497:phyxminiproblems-problemphyxmini0497-incomingalphaspeeddatum, def:physics:phyx-mini-0497:phyxminiproblems-problemphyxmini0497-scatteredalphaspeeddatum, def:physics:phyx-mini-0497:phyxminiproblems-problemphyxmini0497-scatteringangledegreesdatum}
  Numerical source data, the stationary-target boundary condition, and primary image evidence. No field constrains the gold recoil speed
\end{definition}
\begin{proof}
  This is the declared model definition of Matches Problem And Primary Figure.
\end{proof}
\begin{definition}[Satisfies Mass Number Approximation]
  \label{def:physics:phyx-mini-0497:phyxminiproblems-problemphyxmini0497-satisfiesmassnumberapproximation}
  \lean{PhyXMiniProblems.ProblemPhyXMini0497.SatisfiesMassNumberApproximation}
  \uses{def:physics:phyx-mini-0497:phyxminiproblems-problemphyxmini0497-massreadout, def:physics:phyx-mini-0497:phyxminiproblems-problemphyxmini0497-alphagoldscatteringsetup}
  The usual nuclear mass-number approximation: the alpha particle is assigned mass number 4 and the gold isotope mass number 197. Cross multiplication states the same unit-independent mass ratio without dividing by a mass
\end{definition}
\begin{proof}
  This is the declared model definition of Satisfies Mass Number Approximation.
\end{proof}
\begin{definition}[Satisfies Scattering Geometry]
  \label{def:physics:phyx-mini-0497:phyxminiproblems-problemphyxmini0497-satisfiesscatteringgeometry}
  \lean{PhyXMiniProblems.ProblemPhyXMini0497.SatisfiesScatteringGeometry}
  \uses{def:physics:phyx-mini-0497:phyxminiproblems-problemphyxmini0497-velocityreadout, def:physics:phyx-mini-0497:phyxminiproblems-problemphyxmini0497-speedinmeterspersecond, def:physics:phyx-mini-0497:phyxminiproblems-problemphyxmini0497-degreestoradians, def:physics:phyx-mini-0497:phyxminiproblems-problemphyxmini0497-alphagoldscatteringsetup}
  Component form of the geometry shown in the image. The incident alpha defines the positive horizontal axis, and the scattered alpha lies theta above it. This contains no component or magnitude of the gold recoil velocity
\end{definition}
\begin{proof}
  This is the declared model definition of Satisfies Scattering Geometry.
\end{proof}
\begin{definition}[Satisfies Momentum Conservation]
  \label{def:physics:phyx-mini-0497:phyxminiproblems-problemphyxmini0497-satisfiesmomentumconservation}
  \lean{PhyXMiniProblems.ProblemPhyXMini0497.SatisfiesMomentumConservation}
  \uses{def:physics:phyx-mini-0497:phyxminiproblems-problemphyxmini0497-massreadout, def:physics:phyx-mini-0497:phyxminiproblems-problemphyxmini0497-velocityreadout, def:physics:phyx-mini-0497:phyxminiproblems-problemphyxmini0497-alphagoldscatteringsetup}
  Two-dimensional conservation of linear momentum for the isolated alpha--gold interaction. It is stated for every coherent mass and velocity readout. In particular, it determines the recoil vector but does not assume its answer value
\end{definition}
\begin{proof}
  This is the declared model definition of Satisfies Momentum Conservation.
\end{proof}
\begin{lemma}[gold Recoil Speed exact]
  \label{lem:physics:phyx-mini-0497:phyxminiproblems-problemphyxmini0497-goldrecoilspeed-exact}
  \lean{PhyXMiniProblems.ProblemPhyXMini0497.goldRecoilSpeed_exact}
  \uses{def:physics:phyx-mini-0497:phyxminiproblems-problemphyxmini0497-speedinmeterspersecond, def:physics:phyx-mini-0497:phyxminiproblems-problemphyxmini0497-degreestoradians, def:physics:phyx-mini-0497:phyxminiproblems-problemphyxmini0497-alphagoldscatteringsetup, def:physics:phyx-mini-0497:phyxminiproblems-problemphyxmini0497-incomingalphaspeeddatum, def:physics:phyx-mini-0497:phyxminiproblems-problemphyxmini0497-scatteredalphaspeeddatum, def:physics:phyx-mini-0497:phyxminiproblems-problemphyxmini0497-scatteringangledegreesdatum, def:physics:phyx-mini-0497:phyxminiproblems-problemphyxmini0497-matchesproblemandprimaryfigure, def:physics:phyx-mini-0497:phyxminiproblems-problemphyxmini0497-satisfiesmassnumberapproximation, def:physics:phyx-mini-0497:phyxminiproblems-problemphyxmini0497-satisfiesscatteringgeometry, def:physics:phyx-mini-0497:phyxminiproblems-problemphyxmini0497-satisfiesmomentumconservation}
  The exact recoil-speed formula obtained by taking the Euclidean norm of the momentum balance and using the included angle between the two alpha velocities
\end{lemma}
\begin{proof}
  The conclusion follows from the governing relations and figure readouts named in the statement of gold Recoil Speed exact.
\end{proof}
\begin{definition}[Answer Choice]
  \label{def:physics:phyx-mini-0497:phyxminiproblems-problemphyxmini0497-answerchoice}
  \lean{PhyXMiniProblems.ProblemPhyXMini0497.AnswerChoice}
  Labels printed beside the four multiple-choice answers
\end{definition}
\begin{proof}
  This is the declared model definition of Answer Choice.
\end{proof}
\begin{definition}[speed In Meters Per Second]
  \label{def:physics:phyx-mini-0497:phyxminiproblems-problemphyxmini0497-answerchoice-speedinmeterspersecond}
  \lean{PhyXMiniProblems.ProblemPhyXMini0497.AnswerChoice.speedInMetersPerSecond}
  \uses{def:physics:phyx-mini-0497:phyxminiproblems-problemphyxmini0497-speedinmeterspersecond, def:physics:phyx-mini-0497:phyxminiproblems-problemphyxmini0497-answerchoice}
  The recoil speed displayed beside each answer label, in metres per second
\end{definition}
\begin{proof}
  This is the declared model definition of speed In Meters Per Second.
\end{proof}
\begin{definition}[recorded Answer Choice]
  \label{def:physics:phyx-mini-0497:phyxminiproblems-problemphyxmini0497-recordedanswerchoice}
  \lean{PhyXMiniProblems.ProblemPhyXMini0497.recordedAnswerChoice}
  \uses{def:physics:phyx-mini-0497:phyxminiproblems-problemphyxmini0497-answerchoice}
  The answer label recorded in the supplied dataset metadata
\end{definition}
\begin{proof}
  This is the declared model definition of recorded Answer Choice.
\end{proof}
\begin{definition}[Matches Answer Choice]
  \label{def:physics:phyx-mini-0497:phyxminiproblems-problemphyxmini0497-matchesanswerchoice}
  \lean{PhyXMiniProblems.ProblemPhyXMini0497.MatchesAnswerChoice}
  \uses{def:physics:phyx-mini-0497:phyxminiproblems-problemphyxmini0497-planarvelocity, def:physics:phyx-mini-0497:phyxminiproblems-problemphyxmini0497-speedinmeterspersecond, def:physics:phyx-mini-0497:phyxminiproblems-problemphyxmini0497-answerchoice}
  Agreement with a speed displayed to three significant figures. At the scale of answer C, half a unit in the last displayed digit is 500 m/s
\end{definition}
\begin{proof}
  This is the declared model definition of Matches Answer Choice.
\end{proof}
% --- Archon declaration topology end ---
% --- Archon physics formalization source end ---
